\section{Finite search: arithmetic classification and explicit memory bounds}
\label{sec:r80-finite-search-saturation}

The arithmetic obstruction extends to every finite number of search
locations, with arbitrary positive call costs.  The distinction is
between logarithms of the failure probabilities that lie in one
rational direction and those that occupy several rational directions.
The latter case includes families containing both rationally related
and irrationally related pairs.  We obtain a necessary and sufficient
condition for an optimal history law to admit an autonomous stochastic
implementation with finitely many counted labels.  We also give a
primitive-data exponential approximation bound and, for rational data
in the saturating case, a computable bound on the sufficient number of
labels.

The search ordering, interior least-favourable priors, and finite
mixtures of consistent tie orders are classical ingredients
\cite{Matula1964,ClarksonLinGlazebrook2023,BuiLidbetterLin2023}.
Eventual cycles under a common-power relation are also classical for
arbitrary finite $n$; see Section 6.1 and Appendix C of
Clarkson--Lin--Glazebrook~\cite{ClarksonLinGlazebrook2023}, with the
earlier search-theoretic attribution there to Matula.  We give the
required arguments below.  The common-power sufficiency is therefore
a classical search consequence together with counted initialization.
The additional resource conclusions concern all-stochastic necessity
through commensurability classes, explicit primitive-data bounds, and
implementation by counted stochastic memory.  A finite mixture of
infinite words alone is not such an implementation.  This section is
a classification for the displayed search family, not an assertion
that every instrument in Theorem~\ref{thm:r79-quantum-frontier} has an
arithmetic saturation criterion.

Fix $n\ge2$, costs $c_i>0$, and failure probabilities $r_i\in(0,1)$,
for $i\in\{1,\ldots,n\}$.  There are $n$ fixed hidden models and $n$
commands.  In model $i$, command $i$ succeeds with probability
$q_i=1-r_i$, and every other command fails surely.  There is one failure
report, and success stops the experiment.  Command $i$ costs $c_i$ in
every model.  Set
\begin{equation}\label{eq:r80-finite-search-data}
 \tau_i=\frac{c_i}{q_i},\qquad u_i=-\log r_i>0,\qquad
 B=\sum_{i=1}^n\tau_i,\qquad H=\sum_{i=1}^n\frac1{u_i}.
\end{equation}
The hidden model is selected once.  The initialization and memory
conventions are those of Section~\ref{sec:r79-quantum-frontier}:
the initial label is prescribed and pure, and command and update
randomness is permitted only through the fixed controller kernels.
Write $V_M$ and $V=V_{\rm hist}$ for the finite-label and history
values in this family.

\begin{lemma}[A common one-label fallback]
\label{lem:r80-finite-search-one-label}
The one-label value is $V_1=B$.  It is attained by the row
\begin{equation}\label{eq:r80-finite-search-fallback}
 p_i=\frac{q_i^{-1}}{\sum_jq_j^{-1}},
\end{equation}
whose expected cost equals $B$ in every model.
\end{lemma}
\begin{proof}
For a row $p$, put $x_i=p_iq_i$.  A zero $x_i$ gives infinite cost
in model $i$.  Otherwise its one-call mean cost is
$\sum_i p_ic_i=\sum_i\tau_i x_i$, and its worst model cost is
\[
 \frac{\sum_i\tau_i x_i}{\min_i x_i}\ge\sum_i\tau_i=B.
\]
The displayed row makes all $x_i$ equal and attains the bound.
\end{proof}

\subsection{The search saddle and consistent tie orders}

A planned word $w=(w_0,w_1,\ldots)\in\{1,\ldots,n\}^{\mathbb N_0}$
specifies the commands along the hypothetical all-failure path.
Let $k_i(t)$ be the number of occurrences of $i$ before position $t$.
Its cost in model $i$ is the nonnegative sum
\begin{equation}\label{eq:r80-finite-search-word-cost}
 C_i(w)=\sum_{t\ge0}c_{w_t}r_i^{k_i(t)}.
\end{equation}
Every randomized causal policy induces a probability law on these
words; conversely, a word law is realized by its conditional
next-letter kernels.  Indeed, conditional on a command prefix, its
physical failure likelihood is common to all private realizations of
that prefix and cancels.  Thus these conditional kernels do not use
the hidden model.  Tonelli's theorem identifies policy costs with the
expectations of \eqref{eq:r80-finite-search-word-cost}.  Finite cost
in every model requires infinitely many occurrences of every command,
almost surely: after the last occurrence of $i$, its survival factor
is a fixed positive number and every later cost is bounded below by
$\min_jc_j>0$.

\begin{lemma}[The finite-search saddle]
\label{lem:r80-finite-search-saddle}
For an interior prior $\pi=(\pi_1,\ldots,\pi_n)$, put
$a_i=\pi_i/\tau_i$.  A finite-cost word is Bayes optimal if and only if
it chooses, at every prefix, a maximal element of
\begin{equation}\label{eq:r80-finite-search-indices}
 \{a_i r_i^{k_i(t)}:1\le i\le n\}.
\end{equation}
Every such greedy word has finite cost in every model.  If
$b_i=-\log a_i$ and
\[
 R_\pi=\max\{\max_i b_i-\min_i b_i,\max_i u_i\},
 \qquad \bar b=\frac{\sum_i b_i/u_i}{H},
\]
then, uniformly over every choice at ties,
\begin{equation}\label{eq:r80-finite-search-tail-frequency}
 r_i^{k_i(t)}\le e^{b_i-\bar b+R_\pi}e^{-t/H},
 \qquad
 \lim_{t\to\infty}\frac{k_i(t)}t=\frac1{Hu_i}.
\end{equation}

The robust history value is attained.  There is an interior prior
$\pi_*$ with Bayes value $V$, and every robust optimum has cost vector
$(V,\ldots,V)$ and is supported on greedy words for that prior.
One robust optimum is a mixture of at most $n$ greedy words, each of
which resolves every tie using one fixed priority permutation of
$\{1,\ldots,n\}$.
\end{lemma}
\begin{proof}
Exchanging two adjacent different commands $ij$ and $ji$ after counts
$k$ changes only model components $i,j$.  With $ij$ first in the
subtraction the risk difference is
\begin{equation}\label{eq:r80-finite-search-exchange}
 -c_jq_i r_i^{k_i}e_i+c_iq_j r_j^{k_j}e_j,
\end{equation}
where $e_i$ are the coordinate vectors of $\mathbb R^n$.
Its Bayes value is
$c_ic_j(a_jr_j^{k_j}-a_ir_i^{k_i})$.
The greedy rule merges the $n$ decreasing lists
$\{a_i r_i^k:k\ge0\}$ in decreasing numerical order.
Given any finite-cost competitor and a desired greedy next letter,
move the next occurrence of that letter to the desired position.
It exists because every letter occurs infinitely often.  Each
intervening list item is no larger than the desired item: it was no
larger at the starting prefix, and subsequent items in its own list
decrease.  The adjacent exchanges therefore do not increase Bayes
cost, and a strictly incorrect first choice strictly decreases it.
Repeating finitely many times matches any prescribed finite greedy
prefix.  The cost of that prefix is at most the full competitor cost.
Letting its length increase proves optimality of a greedy word once
its finiteness has been verified.  The strict exchange proves
necessity.  Because the prior is interior, any infinite component
cost gives infinite Bayes cost.

For finiteness set $x_i(t)=b_i+u_i k_i(t)$.  A greedy step increments
a current minimum by $u_i$.  Its minimum cannot decrease, and its
maximum is at most the greater of the old maximum and the old minimum
plus $\max_i u_i$.  Induction gives
\[
 \max_i x_i(t)-\min_i x_i(t)\le R_\pi.
\]
The weighted mean with weights $1/u_i$ is
$\bar b+t/H$, since $\sum_i k_i(t)=t$.  It lies between the minimum
and maximum.  Thus $|x_i(t)-\bar b-t/H|\le R_\pi$, which proves
both assertions in \eqref{eq:r80-finite-search-tail-frequency}.
The costs are bounded above by a geometric sum times $\max_i c_i$.

We next prove robust attainment and the interior saddle.  Let
$\mathcal R\subset\mathbb R^n$ consist of the finite upper bounds
on cost vectors attainable by word laws.  This nonempty set is convex
by mixing and upward closed.  It is closed: if finite upper bounds
converge, compactness of probability laws on the finite-alphabet
product space gives a weakly convergent subsequence.  Every
finite-prefix cost is continuous and bounded, so the limiting law
satisfies all the finite-prefix bounds; monotone convergence gives
the full bounds.  The same argument applied to a minimizing sequence
shows robust attainment, and $V\le B$ by
Lemma~\ref{lem:r80-finite-search-one-label}.

Separate $\mathcal R$ from the open orthant
$\{z:z_i<V\text{ for every }i\}$.  Upward closure makes the
nonzero separating vector nonnegative.  Normalize its sum to one,
and call it $\pi_*$.  Since $(V,\ldots,V)\in\mathcal R$, separation
gives $\inf_{z\in\mathcal R}\pi_*\cdot z=V$.
We claim that $\pi_*$ has full support.  Otherwise let
$S=\{i:\pi_{*,i}>0\}$ be a proper subset.  Take an attaining robust
word law and delete every command outside $S$ from every word.
For a model $i\in S$, the survival factor of each retained call is
unchanged, whereas every deleted call had positive cost and strictly
positive survival probability.  Every original word contains all
commands infinitely often almost surely.  Consequently the deleted
law strictly decreases the $\pi_*$-weighted cost, while retaining
finite costs and eventual success for each model in $S$.

The deleted law need not be proper outside $S$.  To obtain admissible
competitors, follow its first $T$ calls and then use the common
fallback \eqref{eq:r80-finite-search-fallback}.  The fallback adds
at most $B$ times the survival probability at $T$ in each model of
$S$.  Those probabilities tend to zero by bounded convergence.
The finite-prefix costs increase to the deleted costs.  These
completed laws are proper in every model and converge in weighted
cost to a value strictly below $V$, a contradiction.  Thus $\pi_*$
is interior.  An attaining robust law has costs at most $V$ in
every component and weighted cost at least $V$, so all its component
costs equal $V$.  Its word law must be supported on Bayes minimizers,
which are precisely the greedy words just characterized.

It remains to prove the finite consistent-tie mixture.  At any tied
numerical level $z>0$, each list contributes at most one item.  The
tied items form a consecutive block; permuting them leaves every
later count unchanged.  For a tied adjacent pair $i,j$, equation
\eqref{eq:r80-finite-search-exchange} becomes
\begin{equation}\label{eq:r80-finite-search-tied-vector}
 c_ic_j z\left(-\frac{e_i}{\pi_{*,i}}
                    +\frac{e_j}{\pi_{*,j}}\right).
\end{equation}
For any linear functional $\lambda\in\mathbb R^n$, its difference
is $c_ic_jz(\lambda_j/\pi_{*,j}-\lambda_i/\pi_{*,i})$.
It follows that the order minimizing $\lambda$ in every tied block
is the same global priority order: decreasing
$\lambda_i/\pi_{*,i}$, with one fixed convention for equal ratios.
Sort finitely many tied blocks in that order and then let their
number increase.  The geometric tail bound
\eqref{eq:r80-finite-search-tail-frequency}, uniform over tie
choices, shows convergence of every component cost.  Hence the
minimum of every linear functional over the greedy risk vectors
is attained by one of the $n!$ consistent priority words.

The greedy word space is closed in the compact product space, and
its risk map is continuous by the same uniform tails.  Its compact
risk image therefore lies in the convex hull of the finitely many
consistent-priority risk vectors: otherwise finite-dimensional
separation would contradict the preceding minimum property.
The robust risk vector, as a barycenter of greedy risks, lies in
that convex hull as well.  All these vectors lie on the hyperplane
$\pi_*\cdot z=V$.  Any convex representation using more than $n$
of them has an affine dependence.  Subtract a suitable multiple
of that dependence from the weights until one weight becomes zero;
the weights remain nonnegative, still sum to one, and give the same
vector.  Repeating leaves at most $n$ vectors, as asserted.
\end{proof}

\subsection{Rigidity of projected finite-state outputs}

\begin{lemma}[A deterministic class suffix from finite stochastic memory]
\label{lem:r80-projected-word-rigidity}
Let an autonomous controller have finitely many labels, finitely many
commands, one failure report, and fixed stochastic command and update
kernels.  Partition the commands into finitely many classes.  If its
planned word law projects, after some fixed finite time, onto one
prescribed class word almost surely, that class word is ultimately
periodic.
\end{lemma}
\begin{proof}
Let $S_t$ be the nonempty set of labels having positive probability
at time $t$ on the hypothetical all-failure path.  The next support is
the union, over labels in $S_t$ and commands of positive probability
in their rows, of the positive supports of the corresponding update
rows.  It is therefore a fixed deterministic function of $S_t$.
At times when the class output is prescribed, every positive command
in every label of $S_t$ belongs to that prescribed class.  Since the
rows have sum one, this class is determined by $S_t$ as well.
There are finitely many nonempty label subsets.  The support-set
trajectory eventually repeats and thereafter is periodic.  Its
prescribed class output is consequently ultimately periodic.  The
argument permits arbitrary stochastic transitions and arbitrary
mixing among commands inside the prescribed class.
\end{proof}

\begin{theorem}[Finite-search saturation for any finite number of locations]
\label{thm:r80-finite-search-saturation}
For the family \eqref{eq:r80-finite-search-data},
\begin{equation}\label{eq:r80-finite-search-iff}
 \bigl(\exists M<\infty:\ V_M=V\bigr)
 \quad\Longleftrightarrow\quad
 \frac{\log r_i}{\log r_j}\in\mathbb Q
       \quad\hbox{for every }i,j.
\end{equation}
If the condition fails, every finite $M$ has the strict attained
gap $V_M>V$, although $V_M\downarrow V$ as $M\to\infty$.
If the condition holds, a robust optimum is realized by a counted
pure-start finite controller, obtained from at most $n$ ultimately
periodic consistent-priority words.
\end{theorem}
\begin{proof}
Partition the indices into classes by $i\sim j$ if $u_i/u_j$ is
rational.  Consider a least-favourable interior prior from
Lemma~\ref{lem:r80-finite-search-saddle}.  Two lists from different
classes have at most one tied pair.  Indeed, two distinct equalities
$a_ir_i^k=a_jr_j^l$ would imply
$(k-k')u_i=(l-l')u_j$ with nonzero integer differences, contrary to
their being in different classes.  There are finitely many pairs
of lists, hence only finitely many cross-class tied levels.  Each
positive numerical level has only finitely many list items above
it.  After a fixed finite position, all remaining ties therefore
occur within a single class.  Ordering the list items then fixes the
same deterministic class suffix for every greedy word, regardless
of the within-class choices.

For a class $C$, \eqref{eq:r80-finite-search-tail-frequency} gives
its positive limiting frequency
\begin{equation}\label{eq:r80-finite-search-class-frequency}
 f_C=\frac{H_C}{H},\qquad H_C=\sum_{i\in C}\frac1{u_i}.
\end{equation}
Choose a representative $u_C$ from each class.  Then
$H_C=A_C/u_C$ for some positive rational $A_C$.
For two distinct classes $C,D$, the ratio
$f_C/f_D=(A_C/A_D)(u_D/u_C)$ is irrational.  An ultimately periodic
class word has rational limiting frequencies, and therefore rational
ratios between any two positive frequencies.  Thus the common class
suffix cannot be ultimately periodic when there is more than one
class.

If a finite controller attained $V$, its word law would be supported
on the greedy words by Lemma~\ref{lem:r80-finite-search-saddle},
so its projected class suffix would be the one just described.
Lemma~\ref{lem:r80-projected-word-rigidity} gives a contradiction.
All failure probabilities in the search family are positive, so
there is no zero-likelihood physical history to remove from this
argument.  The attained-frontier result in
Theorem~\ref{thm:r79-quantum-frontier} makes nonattainment of $V$
equivalent to the strict gap at each finite $M$.  Convergence follows
from Corollary~\ref{cor:r79-history-limit}; an explicit rate is also
proved below.

Conversely suppose that all $u_i/u_j$ are rational.  Write
$u_i=h_i\ell$ with positive integers $h_i$ and $\ell>0$.
For one fixed priority word consider the relative scores
\[
 d_i(t)=x_i(t)-x_n(t),\qquad 1\le i<n.
\]
Each $d_i(t)$ belongs to the fixed coset
$(b_i-b_n)+\ell\mathbb Z$ and lies in
$[-R_{\pi_*},R_{\pi_*}]$.  There are only finitely many such tuples.
The tuple determines the next command by minimizing the scores
relative to $x_n=0$ and using the fixed priority at ties.  On command
$j<n$, only $d_j$ increases, by $u_j$; on command $n$, every $d_i$
decreases by $u_n$.  Thus the tuple evolves deterministically in a
finite set.  The word is ultimately periodic and has a finite
deterministic implementation.

The saddle lemma expresses a robust optimum as a mixture of at most
$n$ such words.  We implement that mixture without a free random
initial label.  Give the implementation one prescribed initial
label and disjoint finite tail controllers for the words.  If their
weights are $\lambda_j$ and first commands are $a_j$, the initial
row chooses command $a$ with probability
$p_a=\sum_{j:a_j=a}\lambda_j$.  For $p_a>0$, after its failure it
selects the corresponding tail $j$ with probability $\lambda_j/p_a$
among those with $a_j=a$.  The update rows for $p_a=0$ are unreachable
and may be assigned arbitrarily.  The physical failure likelihood conditional on $a$
is the same for all those branches in every model, so this realizes
the desired mixture exactly.  The branch identity is stored in the
counted tail labels.  This proves finite saturation.
\end{proof}

\subsection{Bounds from primitive data}

The preceding proof used an interior prior without estimating how
close it can lie to the boundary.  The following bound removes this
unknown from the memory estimates.  Its constants are deliberately
explicit; they are not asserted to be sharp.

\begin{lemma}[An explicit least-favourable prior bound]
\label{lem:r80-finite-search-prior-cap}
Put $c_- =\min_i c_i$, $c_+=\max_i c_i$, $r_- =\min_i r_i$, and
define
\begin{equation}\label{eq:r80-finite-search-prior-constants}
 \begin{aligned}
 T_0&=\left\lceil\frac{2B}{c_-}\right\rceil,
 &g&=\frac{c_-}{2}r_-^{T_0},\\
 T_1&=\left\lceil\frac{2B^2}{c_-g}\right\rceil,
 &K&=c_+T_1+B,\qquad \delta=\frac{g}{2K}.
 \end{aligned}
\end{equation}
Every least-favourable prior whose Bayes value is $V$ satisfies
$\pi_i\ge\delta$ for every $i$.
\end{lemma}
\begin{proof}
Fix an attaining robust word law.  In each model its expected number
of calls is at most $V/c_-\le B/c_-$, since every call costs at
least $c_-$.  Markov's inequality shows that success occurs by call
$T_0$ with probability at least $1/2$.  In model $i$, this implies
that its planned word contains command $i$ among the first $T_0$
positions with probability at least $1/2$.

Delete all occurrences of command $i$.  In any other model $j$, the
survival factor of a retained command is unchanged.  On the event
just described, the first deleted call costs at least $c_-$ and has
model-$j$ survival probability at least $r_-^{T_0}$.  The deleted
law therefore lowers the cost of every model $j\ne i$ by at least
$g$.  Let $W_i$ be the robust history value of the restricted search
problem using only the $n-1$ other models and commands.  We have
\begin{equation}\label{eq:r80-finite-search-deletion-gap}
 W_i\le V-g.
\end{equation}
The same finite-prefix compactness argument used above gives an
attaining restricted law, including when $n-1=1$.

Run this law for at most $T_1$ calls, then switch to the common
fallback of value $B$.  For $j\ne i$, the restricted law has expected
call count at most $W_i/c_-\le B/c_-$.  Its probability of surviving
$T_1$ calls is at most $B/(c_-T_1)$.  The completed policy therefore
has cost at most
\[
 W_i+\frac{B^2}{c_-T_1}\le V-\frac g2
 \quad (j\ne i).
\]
In model $i$ every restricted call fails, but its total cost is at
most $c_+T_1+B=K$.  For a least-favourable prior, this competitor
must have Bayes cost at least $V$.  Hence
\[
 V\le\pi_iK+(1-\pi_i)(V-g/2),\qquad
 \pi_i\ge\frac{g/2}{K-V+g/2}\ge\frac{g}{2K}.
\]
The last inequality uses $V\ge c_->g/2$, which follows from the
cost of the compulsory first call and the definition of $g$.
\end{proof}

\begin{theorem}[An explicit exponential memory approximation]
\label{thm:r80-finite-search-resource-bounds}
With $\delta$ from \eqref{eq:r80-finite-search-prior-constants}, put
\begin{equation}\label{eq:r80-finite-search-rate-constants}
 \begin{aligned}
 \tau_-&=\min_i\tau_i,\qquad \tau_+=\max_i\tau_i,\qquad
 C=\frac{\tau_+}{\delta\tau_-},\\
 A&=C\max\{C,r_-^{-1}\},\qquad
 \kappa=e^{-1/H}\in(0,1).
 \end{aligned}
\end{equation}
For every $M\ge n$,
\begin{equation}\label{eq:r80-finite-search-resource-rate}
 0\le V_M-V\le BA\,\kappa^{\lfloor M/n\rfloor}.
\end{equation}
The bound holds whether or not finite saturation occurs.
\end{theorem}
\begin{proof}
Choose the optimal mixture of at most $n$ consistent-priority words
from Lemma~\ref{lem:r80-finite-search-saddle}.  By the prior bound,
$a_i=\pi_{*,i}/\tau_i$ lies between $\delta/\tau_+$ and
$1/\tau_-$.  Put $L=\log C$ and $R=\max\{L,\max_i u_i\}$.
Then the initial logarithmic-score spread is at most $L$, and
$R_{\pi_*}\le R$.  Also $b_i-\bar b\le L$ and
$e^{R+L}=C\max\{C,r_-^{-1}\}=A$.
Equation~\eqref{eq:r80-finite-search-tail-frequency} now bounds
every model's survival probability after $T$ commands, in each word
and hence in the mixture, by $A\kappa^T$.

Retain the first $T\ge1$ commands of each mixture word and then use
the one-label fallback.  The tail replacement increases each cost
by at most $BA\kappa^T$.  A counted implementation uses one initial
label, at most $n(T-1)$ disjoint continuation labels, and one common
fallback label.  The initial branching is implemented after the
first command exactly as in the saturation theorem.  Since $n\ge2$,
the total is at most $2+n(T-1)\le nT$.  Taking
$T=\lfloor M/n\rfloor$ and adding unused labels if necessary proves
the claim.
\end{proof}

\begin{corollary}[A decidable rational criterion and a sufficient label cap]
\label{cor:r80-finite-search-effective}
Suppose that $r_i,c_i$ are positive rational numbers with $r_i<1$.
For each $i$, let $\nu(r_i)$ be its signed vector of prime exponents,
using the finite union of primes occurring in all numerators and
denominators.  Finite saturation occurs if and only if all these
vectors lie on the same positive rational ray.  This is decidable
by factorization and integer comparisons.

In the saturating case one can compute $t\in\mathbb Q\cap(0,1)$
and positive integers $h_i$ with $r_i=t^{h_i}$.  Using the rational
constants $\delta,\tau_-,\tau_+$ above, choose an integer $J\ge0$
such that
\[
 \begin{aligned}
 t^{-J}&\ge\frac{\tau_+}{\delta\tau_-},
 \qquad N_0=\max\{J,\max_i h_i\},\\
 N&=(2N_0+1)^{n-1},\qquad M_0=1+nN.
 \end{aligned}
\]
Then $V_{M_0}=V$.  In particular, the history value in this
saturating rational family is a real algebraic number and is
computable by the finite quantified description of
Theorem~\ref{thm:r79-quantum-frontier} at the explicitly bounded
budget $M_0$.
\end{corollary}
\begin{proof}
For rational numbers, a rational relation between their logarithms
is equivalent to a positive integer power identity.  Unique prime
factorization makes this equivalent to positive rational
collinearity of their signed exponent vectors.  This proves the
criterion using Theorem~\ref{thm:r80-finite-search-saturation}.
To construct the common base, divide $\nu(r_1)$ by the greatest
common divisor of its nonzero coordinates in absolute value, giving
a primitive integer vector $w$.  Each other exponent vector is
$h_iw$ with $h_i>0$.  Since the coordinates of $w$ have greatest
common divisor one, an integer linear combination of them is one;
applying it to $h_iw$ shows that $h_i$ is an integer.  Then
$t=\prod_p p^{w_p}$ is rational, satisfies $0<t<1$, and has
$r_i=t^{h_i}$.  A common divisor of all $h_i$ may be absorbed in
$t$ if desired.

All constants in \eqref{eq:r80-finite-search-prior-constants} are
rational and effectively obtained by finite arithmetic and ceilings.
The inequality defining $J$ is found by increasing $J$ and making
exact rational comparisons.  Write $\ell=-\log t$.  It implies
$L=\log C\le J\ell$, while $u_i=h_i\ell$, so
$R=\max\{L,\max_i u_i\}\le N_0\ell$.
Each relative-score coordinate in the saturation proof belongs to
one fixed coset of $\ell\mathbb Z$ inside
$[-N_0\ell,N_0\ell]$.  It has at most $2N_0+1$ possible values,
independently of the unknown offset of that coset.  Each consistent
priority word therefore uses at most $N$ relative-score states.
The counted pure-start mixture uses at most one initial label and
$n$ copies of those state sets, giving $M_0=1+nN$.

The family has rational scalar instrument and cost data.  Its
finite-budget optimum is attained, finite, algebraic, and specified
by the real-algebraic quantified formula of
Theorem~\ref{thm:r79-quantum-frontier}.  Since the proven cap gives
$V=V_{M_0}$, that finite procedure computes the history value here.
The procedure is an effectivity result; no practical complexity
bound for quantifier elimination is asserted.
\end{proof}

The two-location normalized bounds in
Section~\ref{sec:r80-arithmetic-saturation} retain their sharper
constants.  The present results add arbitrary finite model number,
arbitrary positive deterministic command costs, mixed
commensurability classes, a primitive-data convergence rate, and an
effective sufficient memory cap in the rational saturating case.
They do not infer a positive numerical lower bound on the gap
$V_M-V$ from irrationality alone, nor a comparable arithmetic
classification for general noncommuting instruments.
